\subsection{Reflection and the exact integral torsion formula}
\label{ir:sec:reflection}

\subsubsection{Reflection is a separate quotient operation}
The map $c(e_x)=e_{-x}$ preserves the relations and satisfies $c^2=1$.
It acts on the entire resolution by $c[S,x]=[S,-x]$.
For integer weights $a_p$, let $D=\mathcal D_q(\Z,a)$ and define
\[
 D^+=D/(c-1)D,\qquad D^-=D/(c+1)D.
\]
These are \emph{coinvariants with a prescribed sign}, not the submodules
of invariant and anti-invariant elements. Over a field of characteristic
not two, the corresponding eigenspaces and coinvariants agree. Over
$\Z$ this distinction is essential.

For any $q$, put
\begin{equation}\label{ir:eq:r-h}
 r(q)=\#\{p>2:p\mid q\}+\mathbf 1_{4\mid q}.
\end{equation}
For $q>2$, we have $r(q)\ge1$ and set
$n(q)=\varphi(q)/2$, $h(q)=2^{r(q)-1}$.

\begin{theorem}[Weighted integral reflection theorem]
\label{ir:thm:integer-reflection}
Let $q>2$ and specialize every $A_p$ to an integer $a_p$. For each
$\eps\in\{1,-1\}$,
\begin{equation}\label{ir:eq:integer-reflection}
 D/(c-\eps)D\simeq\Z^{n(q)}\oplus(\Z/2\Z)^{t(q,a)},
\end{equation}
where
\begin{equation}\label{ir:eq:torsion-count}
 t(q,a)=
 \begin{cases}
  h(q),&a_p\text{ is odd for every odd prime }p\mid q,\\
  0,&\text{some odd prime }p\mid q\text{ has }a_p\text{ even}.
 \end{cases}
\end{equation}
In particular the answer is independent of the integer $a_2$.
At $q=1$ or $q=2$, the separate answers are $D^+=\Z$ and
$D^-=\Z/2\Z$.
\end{theorem}

The absence of odd torsion is an elementary property of an involution
on a free lattice. The exact multiplicity is the substantive calculation.
We prove it through the fixed points of the resolution.

\subsubsection{Two elementary facts about involutions}
For a module $M$ with an involution $c$, define its two periodic
cohomology groups by
\begin{equation}\label{ir:eq:tate-definition}
 \widehat H^0(M)=\ker(c-1)/(1+c)M,\qquad
 \widehat H^1(M)=\ker(1+c)/(c-1)M,
\end{equation}
and repeat with period two. These are the cohomology groups of the
complex whose differentials alternate between $c-1$ and $1+c$,
with $c-1$ leaving even degree. This definition, rather than any
unmentioned cohomology convention, fixes all signs and parity labels.

\begin{lemma}[Torsion in sign coinvariants]\label{ir:lem:torsion-tate}
If $M$ is a finite-rank free abelian group, every torsion class in
$M/(c-\eps)M$ is killed by $2$. Its torsion subgroup is
$\widehat H^1(M)$ for $\eps=1$, and $\widehat H^0(M)$ for $\eps=-1$.
\end{lemma}
\begin{proof}
Suppose $k v=(c-\eps)w$ with $k\ne0$. Applying $c+\eps$ gives
$k(c+\eps)v=0$, hence $(c+\eps)v=0$ because $M$ is torsion-free.
Then $(c-\eps)v=-2\eps v$, so the class of $v$ is killed by $2$.
Conversely a vector in $\ker(c+\eps)$ represents a class killed by
$2$. Taking its quotient by $(c-\eps)M$ gives exactly the two groups
in \eqref{ir:eq:tate-definition}.
\end{proof}

\begin{lemma}[Balanced free ranks]\label{ir:lem:balanced}
For every field of characteristic different from two, and arbitrary
weights in that field, the two reflection coinvariants of $\mathcal D_q$
have dimension $\varphi(q)/2$ when $q>2$.
\end{lemma}
\begin{proof}
On a permutation module with $u$ basis points and $f$ fixed points,
the two eigenspace dimensions are the integers $(u+f)/2$ and
$(u-f)/2$. The sign-eigenspace functors are exact when $2$ is
invertible, so apply this count to the resolution. The alternating
sum of the $u$'s is $\varphi(q)$ by \eqref{ir:eq:euler-rank}.
The fixed-point count for $X_m$ is one if $m$ is odd and two if
$m$ is even. Thus its alternating sum is
\[
 \sum_{S\subseteq P(q)}(-1)^{|S|}\gcd(2,q/p_S)=0
 \qquad(q>2).
\]
For odd $q$, this is $(1-1)^{|P(q)|}$. If $v_2(q)=1$, summing
first over whether $2\in S$ leaves $(2-1)(1-1)^k$, where
$k$ is the number of odd prime divisors and is positive. If
$v_2(q)\ge2$, that first sum is $2-2$. Hence both dimensions
are $\varphi(q)/2$.

This uses integer dimension counts in the chain modules, not merely
a trace equality in the coefficient field. The latter alone would
not determine dimensions in positive characteristic.
\end{proof}

\subsubsection{Removing every nonfixed orbit from periodic cohomology}
This step explains why the calculation is small despite the $q$
original grid points. Let $C^{\mathrm{mov}}_j$ be the span of the
nonfixed points in each summand of $C_j$. It is a subcomplex:
if $x\ne-x$ and $py=x$, then $y$ cannot be fixed; the term
$-a_p[S\setminus\{p\},x]$ is also nonfixed. Each term of this
subcomplex is a direct sum of copies of the regular permutation
module $S[\Z/2\Z]$ over the chosen coefficient ring $S$.

The periodic complex of a regular two-point orbit is exact over
\emph{every} commutative ring. Indeed the kernel of the difference
map is the diagonal, which is the image of the sum map; the kernel
of the sum map consists of $(u,-u)$, which is the image of the
difference map. This assertion requires neither division by $2$
nor absence of $2$-torsion.

Let $K_\bullet=C_\bullet/C^{\mathrm{mov}}_\bullet$. Its generators
are only the fixed symbols $[S,0]$ and, when available, $[S,1/2]$.
The action on $K_\bullet$ is trivial. Forming the double complex of
the periodic cohomology complexes gives
\begin{equation}\label{ir:eq:tate-comparison}
 \widehat H^\bullet(\mathcal D_q(S,a))\simeq
 H^\bullet\operatorname{Tot}\bigl(T(K_\bullet)\bigr).
\end{equation}
Here $T$ denotes the alternating difference/sum complex just defined,
and the chain degree of $K_j$ contributes $-j$ to total degree.

For completeness, this comparison can be justified without assuming a
spectral sequence degenerates. The augmented resolution is exact after
base change, so taking horizontal homology in the double complex
replaces $C_\bullet$ by $D$ in degree zero. Removing the moving
subcomplex does not change total cohomology because each of its
vertical columns is exact. In both statements the horizontal width
is bounded by $|P(q)|$, so in each total degree only finitely many
columns occur; successive elimination along these columns justifies
the two comparisons. In particular, no convergence of an infinite
horizontal filtration is being used.

\subsubsection{The fixed-point complex is a Koszul complex}
For a sequence $b_1,\ldots,b_r$ in a ring $S$, let
$K_S(b_1,\ldots,b_r)$ denote the tensor product of the two-term
complexes $S\xrightarrow{b_i}S$, with source in chain degree one.
This is the Koszul complex; its definition is all that is needed below.

For an odd prime $p$, the fixed-point part of its distribution map
is multiplication by $1-a_p$ on each available fixed symbol. Of the
$p$ roots of a fixed point, exactly one is fixed, and is the original
point. The other roots disappear in the quotient by moving orbits.
The factor at $2$ has two different forms.

When $v_2(q)=1$, its source has only $0$ and its target has $0,1/2$.
Its column is
\begin{equation}\label{ir:eq:two-column}
 \begin{pmatrix}1-a_2\\1\end{pmatrix}.
\end{equation}
The unit entry splits off an identity complex, leaving one free
copy of $S$ in degree zero. Therefore this prime contributes no
Koszul generator to the fixed-point complex.

When $v_2(q)\ge2$, both grids have two fixed points. In the order
$(0,1/2)$, the matrix is
\begin{equation}\label{ir:eq:two-matrix}
 B(a_2)=\begin{pmatrix}1-a_2&0\\1&-a_2\end{pmatrix}.
\end{equation}
The two roots of $1/2$ are $1/4,3/4$, which are moving points; this
explains the second column. Unimodular row and column operations
transform this matrix to $\operatorname{diag}(1,a_2(a_2-1))$, up to
a unit sign. These operations are valid over $\Z[a_2]$ and over
every specialization, because they pivot on the displayed $1$.
Odd-prime maps are scalar, so the operations commute with all other
directions of the complex.

We have proved the following ring-uniform reduction.
\begin{proposition}[Fixed-point reduction]\label{ir:prop:fixed-koszul}
Over any commutative ring, $K_\bullet$ is chain-homotopy equivalent to
\begin{equation}\label{ir:eq:active-koszul}
 K_S(\boldsymbol\delta),\qquad
 \boldsymbol\delta=
 \bigl(a_p-1:\ p>2,\ p\mid q\bigr)
 \ \text{together with }a_2(a_2-1)\text{ if }4\mid q.
\end{equation}
It has $r(q)$ active scalar entries. At $q=1,2$, this is the empty
Koszul complex $S$ in degree zero.
\end{proposition}

\subsubsection{Completing the integer calculation}
For $S=\Z$, the periodic complex on a trivial module alternates
between $0$ and multiplication by $2$. It separates into adjacent
two-term blocks. Tensoring one such block with the free complex
$K_\Z(\boldsymbol\delta)$ computes
$K_\Z(\boldsymbol\delta)\otimes\mathbb F_2$: multiplication by $2$ is a
free resolution of reduction modulo two. Consequently
\begin{equation}\label{ir:eq:integer-tate-koszul}
 \widehat H^\eps(D)\simeq
 \bigoplus_{\substack{0\le j\le r(q)\\j\equiv\eps\ (2)}}
 H_j\bigl(K_{\mathbb F_2}(\overline{\boldsymbol\delta})\bigr),
 \qquad\eps=0,1.
\end{equation}
Only an isomorphism of abelian groups is asserted here; no extra
multiplicative structure is required.

If some odd $a_p$ is even, its active scalar reduces to $1$, so its
two-term complex is contractible and all the homology vanishes.
Otherwise every odd-prime scalar is zero modulo two. The remaining
scalar $a_2(a_2-1)$, when present, is also zero modulo two for every
integer $a_2$. All differentials in the reduced active complex
therefore vanish, and
\[
 \dim_{\mathbb F_2}H_j=\binom{r(q)}j.
\]
For $r(q)\ge1$, the sums over even and odd $j$ are both
$2^{r(q)-1}$. Lemmas~\ref{ir:lem:torsion-tate} and
\ref{ir:lem:balanced} now prove
Theorem~\ref{ir:thm:integer-reflection}. When $r(q)=0$, the only
homology is in degree zero, giving the stated distinct answers for
$q=1,2$.

\subsubsection{Examples and the classical specialization}
At level four, put $a=a_2$. The basis is $(e_0,e_{3/4})$, and
\[
 e_{1/2}=(a-1)e_0,\qquad e_{1/4}=a(a-1)e_0-e_{3/4},
 \qquad
 c=\begin{pmatrix}1&a(a-1)\\0&-1\end{pmatrix}.
\]
For integer $a$, the nonzero Smith factor of either $c-1$ or $c+1$
is $2$, because $a(a-1)$ is even. This two-dimensional example
exhibits both the unavoidable parity and the irrelevance of the
integer value of $a_2$.

For all odd-prime weights odd, the following values result:
\[
\begin{array}{c|rrrrrr}
 q&4&12&15&30&60&105\\\hline
 n(q)&1&2&4&4&8&24\\
 r(q)&1&2&2&2&3&3\\
 h(q)&1&2&2&2&4&4
\end{array}
\]
For example, either sign at level $60$ gives
$\Z^8\oplus(\Z/2)^4$ in the exceptional parity case, but only
$\Z^8$ as soon as $a_3$ or $a_5$ is even.

When all $a_p=1$, the count agrees with the classical ordinary
sign-cohomology theorem of Kubert \cite{intdist:KubertSign}; twice an odd
level has the same effective prime count as that odd level. The
weighted parity criterion, the fixed-point matrices, and the later
jet calculation are the present explicit continuation of the
manuscript's integral question.
